\documentclass{beamer}

\usetheme{Madrid}
\usecolortheme{default}

\usepackage[utf8]{inputenc}
\usepackage[T1]{fontenc}
\usepackage{amsmath}
\usepackage{graphicx}
\usepackage{hyperref}
\title{Model Straussa - sieć o zadanej liczbie krawędzi i trójkątów}
\author{Jan Biały}
\institute{Politechnika Warszawska}
\date{\today}

\begin{document}

%------------------------------------------------

\begin{frame}
\titlepage
\end{frame}

%------------------------------------------------

\begin{frame}{Czym jest model Straussa?}

Model Straussa jest probabilistycznym modelem sieci złożonych, zaproponowanym początkowo do opisu sieci społecznych.

\vspace{0.5cm}

Główne założenia:

\begin{itemize}
    \item sieć traktowana jest jako graf losowy,
    \item prawdopodobieństwo wystąpienia grafu zależy od jego struktury,
    \item model kontroluje:
    \begin{itemize}
        \item liczbę krawędzi,
        \item liczbę trójkątów.
    \end{itemize}
\end{itemize}

\vspace{0.4cm}

Model należy do klasy \textbf{wykładniczych grafów losowych (ERG - Exponentially Random Graph)}.

\end{frame}

%------------------------------------------------

\begin{frame}{Matematyczny opis modelu}

Każda siec ma określone prawdopodobieństwo realizacji $P(G)$. W wykładniczych grafach przypdakowych jest ono opisane funkcją wykładniczą:

\[
P(G)\propto e^{H(G)}
\]

gdzie Hamiltonian sieci ma postać:

\[
H(G)=\theta E(G)+\sigma T(G)
\]

Parametry:

\begin{itemize}
    \item $E(G)$ -- liczba krawędzi,
    \item $T(G)$ -- liczba trójkątów,
    \item $\theta$ -- kontroluje powstawanie krawędzi,
    \item $\sigma$ -- kontroluje tendencję do tworzenia skupień.
\end{itemize}

\end{frame}

%------------------------------------------------

\begin{frame}{Dlaczego trójkąty są ważne?}

Trójkąty opisują lokalne skupiska połączeń w sieci.

\vspace{0.5cm}

W bardzo dużym skrócie można to wytłumaczyć w ten sposób:
\begin{itemize}
    \item Osoba A zna osobę B
    \item Osoba B zna osobę C
    \item Rośnie prawdopodobieństwo, że A zna C
\end{itemize}

\vspace{0.5cm}

Wpływ parametru $\sigma$:

\begin{itemize}
    \item małe $\sigma$ $\rightarrow$ sieć rzadka,
    \item duże $\sigma$ $\rightarrow$ sieć gęsta, silne klastrowanie.
\end{itemize}

\end{frame}

%------------------------------------------------

\begin{frame}{Symulacja Monte Carlo i przejścia fazowe}

Dokładne rozwiązanie modelu jest trudne analitycznie.

\vspace{0.4cm}

Dlatego wykorzystuje się metodę Monte Carlo:

\begin{enumerate}
    \item losujemy parę wierzchołków,
    \item proponujemy dodanie lub usunięcie krawędzi,
    \item obliczamy zmianę energii:
\end{enumerate}

\[
\Delta H = H(G_j)-H(G_i)
\]

Akceptacja zmiany następuje zgodnie z kryterium probabilistycznym.

\vspace{0.4cm}

W modelu obserwujemy:

\begin{itemize}
    \item fazę rzadką,
    \item fazę gęstą,
    \item nagłe przejścia fazowe.
\end{itemize}

\end{frame}

%------------------------------------------------

\begin{frame}{Histereza i zastosowania}

Model Straussa wykazuje zjawisko:

\[
\textbf{Histerezy}
\]

czyli zależności aktualnego stanu układu od jego historii.

\vspace{0.5cm}

Oznacza to możliwość występowania nieciągłych przejść fazowych.

\vspace{0.4cm}

Przykładowe zastosowania:

\begin{itemize}
    \item analiza sieci społecznych,
    \item modelowanie powstawania społeczności,
    \item sieci biologiczne,
    \item badania z zakresu fizyki statystycznej.
\end{itemize}

\vspace{0.4cm}

Cel projektu: implementacja oraz analiza numeryczna modelu Straussa.

\end{frame}

%------------------------------------------------

%------------------------------------------------

%------------------------------------------------

\begin{frame}{Źródła}

\footnotesize

\begin{itemize}

\item Park J., Newman M.E.J.\\
\textit{Solution for the properties of a clustered network}\\
Phys. Rev. E 72 (2005)

\vspace{0.3cm}

\item Fronczak A.\\
\textit{Wykładnicze grafy przypadkowe: teoria, przykłady, symulacje numeryczne}\\
\url{https://if.pw.edu.pl/~agatka/pub/2014rozdzERG.pdf}

\vspace{0.3cm}

\item Robins G., Pattison P., Kalish Y., Lusher D.\\
\textit{Introduction to exponential random graph models}\\
\url{https://doi.org/10.1016/j.socnet.2006.08.002}

\vspace{0.3cm}

\item Exponential Random Graph Models (ERGM)\\
\url{https://en.wikipedia.org/wiki/Exponential_family_random_graph_models}

\end{itemize}

\end{frame}

\end{document}